%%%%%%%%%%%%%%%%%%%%%%%%%%%%%%%%%%%%%%%%%%%%%%%%%%%%%%%%%%%%%%
%% Chapter 2: Related Work
%%%%%%%%%%%%%%%%%%%%%%%%%%%%%%%%%%%%%%%%%%%%%%%%%%%%%%%%%%%%%%

\section{Related Work}
\label{chap:related-work}

\subsection{Sign language recognition and translation (SLR/SLT)}
\label{sec:slr-slt-background}

Automatic sign language translation decomposes around the intermediate,
gloss-level representation introduced in Chapter~\ref{chap:introduction}
(Section~\ref{sec:motivation-context}). This decomposition yields four directed
sub-tasks: \emph{video-to-gloss} and \emph{gloss-to-text} on the recognition
side, and \emph{text-to-gloss} and avatar-based \emph{gloss-to-video} on the
production side. The present project targets the video-to-gloss stage first
(Chapter~\ref{chap:isolated-gloss-recognition}), then extends toward gloss-text
translation (Chapter~\ref{chap:text-translation}) and their integration into a
sentence-level pipeline (Chapter~\ref{chap:sentence-level-translation}).

Even in its isolated form -- where each input video contains a single,
pre-segmented sign -- video-to-gloss recognition is non-trivial. The same gloss
is executed differently across signers, varying in hand shape, movement
trajectory, signing pace, and body proportions, producing substantial
intra-class variability. Conventional softmax classifiers, which partition the
input space with fixed decision boundaries, are poorly suited to absorbing this
variability, particularly under the limited-data conditions typical of sign
language corpora. This observation motivates the retrieval-based, metric-learning
formulation adopted for isolated gloss recognition in this work (see
Chapter~\ref{chap:isolated-gloss-recognition}), in contrast to the prevailing
classification-based approaches surveyed below.

\subsection{Video-to-gloss methods}
\label{sec:video-to-gloss-methods}

This section surveys architectures proposed in the literature for isolated and
continuous sign gesture recognition, spanning skeleton-based, appearance-based,
and hybrid approaches.

\emph{Sign Pose-Based Transformer for Word-Level Sign Language Recognition}
\cite{bohacek2022spoter} introduces SPOTER, a system that discards raw RGB frames
entirely and operates only on 2D skeletal keypoints extracted from a pose
estimator (body and hand joints). Each frame is represented as a vector of joint
coordinates; the sequence of these vectors across the video clip forms the input
to a standard transformer encoder-decoder, where a single learnable class query
token attends over the encoded frame sequence (similar to the DETR
object-detection paradigm) to produce a final classification. Key contributions
include a normalization step that rescales joint coordinates relative to a fixed
body reference frame to remove subject- and camera-distance variance, and a
skeleton-specific augmentation scheme (random joint rotation, scaling, and
perspective-like distortion applied directly to keypoints rather than pixels).
Reported results show 63.18\% top-1 accuracy on the larger, multi-signer WLASL
dataset versus 100\% on the smaller, more controlled LSA64 dataset, indicating
that the lightweight architecture trades some accuracy on diverse, large-vocabulary
data for a substantial reduction in model size and inference cost.

\emph{Power Data Classification: A Hybrid of a Novel Local Time Warping and LSTM}
\cite{li2016ltwlstm} follows a standard LSTM formulation for sequence
classification, with softmax over the summed hidden states, but its distinguishing
contribution is Local Time Warping (LTW), a custom distance metric that replaces
the dynamic-programming alignment of DTW with a non-commutative, fixed-window
local comparison -- reducing complexity from $O(nm)$ to linear time while
retaining comparable discriminative power. The final classifier fuses a
1-nearest-neighbour classifier using LTW with the LSTM's softmax output at the
probability level. Across five-fold cross-validation, the standalone LSTM reaches
roughly 85--88\% accuracy, comparable to the LTW-based nearest-neighbour classifier
alone, but fusing the two branches raises accuracy to 89--93\%, a result the
authors attribute to the parametric and non-parametric branches making
structurally different errors. This dual-branch fusion pattern is structurally
analogous to combining a learned temporal encoder with FAISS-based retrieval, one
of the extensions discussed in Section~\ref{sec:gloss-recognition-discussion}.

\emph{A Comprehensive Study on Deep Learning-Based Methods for Sign Language
Recognition} \cite{adaloglou2021comprehensive} re-implements and retrains several
model families under identical experimental conditions (same datasets,
preprocessing, and splits) to isolate architecture effects from confounding
factors that typically vary across papers. The evaluated families include 2D-CNN
frame classifiers, 3D-CNNs (e.g.\ I3D-style spatiotemporal convolutions),
CNN+RNN hybrids, and graph-based models operating on pose keypoints, evaluated
separately for isolated sign recognition and continuous sign language recognition
(CSLR, typically trained with a CTC loss). The central finding is that no single
architecture dominates across all datasets and tasks: pose-based and graph-based
models tend to generalise better across signers, while appearance-based CNN/3D-CNN
models tend to score higher on datasets with limited signer diversity --
underscoring that reported accuracy figures are heavily dataset-dependent. This
methodological point directly motivates the controlled, single-condition
benchmark of five temporal encoders conducted in
Chapter~\ref{chap:isolated-gloss-recognition}.

A broader review of architectures across both recognition and translation
\cite{review2024dlprocessing} covers 3D-CNN approaches that apply spatiotemporal
convolution kernels directly over stacked video frames, and Transformer-based
sign language translation systems (e.g.\ Camgoz et al.) that replace the CNN+RNN
pipeline with self-attention over either video features or extracted gloss
sequences, removing the sequential bottleneck of RNNs. It also surveys
multimodal transfer-learning baselines that pretrain visual encoders on
large-scale action-recognition corpora (e.g.\ Kinetics) before fine-tuning on
smaller sign language datasets, consistently improving downstream recognition
accuracy relative to training from scratch.

Several further studies characterise the practical range of appearance- and
hybrid-based approaches. A GoogLeNet-plus-BiLSTM pipeline trained on
field-recorded (uncontrolled illumination and background) video
\cite{agri2021bilstm} reaches 76.21\% average accuracy on a custom
agricultural-domain dataset, notably lower than figures reported by comparable
hybrid models trained on studio-recorded data -- attributed by the authors to the
uncontrolled acquisition setting rather than to a weakness of the architecture. A
comparative benchmark of CNN, MLP, RNN, SVM, and KNN classifiers on static
hand-shape images \cite{ann2022recognition} reaches 99\% effectiveness on a
50{,}000-image training set, though this reflects training-set fit rather than
generalisation, and the underlying model does not handle continuous or dynamic
signs. A modality-level review \cite{modalities2022review} separates
vision-based (RGB/RGB-D feeding CNN or CNN-RNN pipelines), skeleton/pose-based
(keypoint sequences feeding GCNs or transformers), and wearable-sensor-based
(IMU/EMG feeding RNNs) systems, noting that multimodal benchmarks such as AUTSL
generally report higher accuracy than single-modality RGB datasets due to the
added discriminative information from depth and skeletal streams.

Two further architectures combine convolutional and recurrent stages directly.
A 3D-CNN front end stacked with LSTM layers \cite{cnn3dlstm2023realtime}
convolves across height, width, and time simultaneously to extract motion-aware
spatial features without a separate optical-flow stage, before passing the
resulting per-clip feature sequence to LSTM layers for longer-range temporal
modelling. \emph{Deepsign} \cite{deepsign2022} stacks a single LSTM layer
followed by a GRU layer over frame-level features, and empirically finds that,
of the four possible two-layer LSTM/GRU orderings, LSTM-then-GRU performs best,
reaching approximately 97\% accuracy across eleven gesture classes on a custom
dataset.

Finally, \emph{Spatial-Temporal Graph Convolutional Networks for Sign Language
Recognition} \cite{correia2019stgcn} represents each frame's skeletal pose as a
graph, with joints as nodes and anatomical connections (e.g.\ wrist--elbow,
elbow--shoulder) as edges. A spatial graph convolution aggregates information
between connected joints within a frame, while a temporal convolution connects
the same joint across consecutive frames, following the ST-GCN formulation
originally developed for general skeleton-based action recognition. This
explicitly encodes kinematic structure as an inductive bias, in contrast to
approaches that receive only a flattened coordinate vector per frame -- a
limitation of the model architecture adopted in
Chapter~\ref{chap:isolated-gloss-recognition} and revisited in
Section~\ref{sec:gloss-recognition-discussion} as a direction for future work.

\subsection{NLP models for sign language}
\label{sec:nlp-models-sign-language}

Beyond visual recognition, gloss-to-text and text-to-gloss translation cast sign
language processing as a sequence-to-sequence NLP problem, typically approached
with Transformer-based encoder-decoder architectures operating on gloss token
sequences rather than raw text alone (see the Transformer-based translation
systems noted in Section~\ref{sec:video-to-gloss-methods}
\cite{review2024dlprocessing}).

\todo{PLACEHOLDER: Survey NLP models specifically applied to gloss2text /
text2gloss translation -- e.g. sequence-to-sequence and Transformer-based sign
language translation models (Camgoz et al., Sign Language Transformers),
pretrained language model fine-tuning for gloss translation, back-translation
and data augmentation strategies for low-resource gloss corpora. This section
should ground the modelling choices made in Chapter~\ref{chap:text-translation}.}

\subsection{Existing systems and benchmarks}
\label{sec:existing-systems-benchmarks}

The literature surveyed above relies on a range of datasets that also inform the
benchmark design in this thesis. \textbf{WLASL} and \textbf{LSA64} are used in
\cite{bohacek2022spoter} to contrast performance on a large, multi-signer
vocabulary against a smaller, controlled one. \textbf{AUTSL} is a large-scale
multimodal Turkish sign language dataset used to benchmark CNN-LSTM and 3D-CNN
baselines \cite{modalities2022review}. \textbf{ASLLVD} underlies the
skeleton-based dataset released alongside the ST-GCN model
\cite{correia2019stgcn}. \textbf{IISL2020} is a custom eleven-class dataset used
to benchmark the stacked LSTM-GRU architecture in \cite{deepsign2022}. The
present work benchmarks isolated gloss recognition on a subset of the
\textbf{ASL Citizen} dataset (Microsoft, 2023), detailed in
Section~\ref{sec:gloss-recognition-dataset}.

\todo{PLACEHOLDER: Describe existing end-to-end SLR/SLT systems and public
benchmarks/leaderboards (e.g. How2Sign, RWTH-PHOENIX-Weather 2014T, existing
production or research demo systems) beyond the datasets already covered above,
and relate them to the scope of this thesis.}
